% Cue conflict in chain-of-thought arithmetic -- ARR short-paper draft (4 pages of body).
%
% Every quantity is a \NUM{} macro from paper/numbers.tex, written by scripts/51_tables.py from
% results/summary/; an undefined key renders in red. No number in the prose is typed by hand.
%
% Build: make paper   (tectonic at envs/tex/bin/tectonic; prints body page count vs limit 4)

\documentclass[11pt]{article}
\usepackage[review]{acl}      % anonymous + line numbers for ARR; final / preprint later
\usepackage{times}
\usepackage{latexsym}
\usepackage[T1]{fontenc}
\usepackage[utf8]{inputenc}
\usepackage{microtype}
\usepackage{inconsolata}
\usepackage{graphicx}
\usepackage{booktabs}
\usepackage{amsmath}
\usepackage{xcolor}
\usepackage{url}

% An unfilled number must be impossible to miss.  % grep-skip
\providecommand{\NUM}[1]{\textcolor{red}{$\langle\langle$\texttt{#1}$\rangle\rangle$}}
\providecommand{\tabinput}[1]{\IfFileExists{tables/#1.tex}{\input{tables/#1.tex}}{\textcolor{red}{[table \texttt{#1} not generated yet]}}}
\IfFileExists{numbers.tex}{\input{numbers.tex}}{}

\definecolor{neu}{RGB}{60,60,60}
\definecolor{con}{RGB}{0,110,60}
\definecolor{inc}{RGB}{180,30,30}

\title{Chain of Thought Reduces Errors from Misleading Names in Arithmetic}

\author{Anonymous ACL submission}

\begin{document}
\maketitle

\begin{abstract}
Asking language models to write out arithmetic calculations reduces answers that follow
a misleading variable name. Building on \citeauthor{kudo2026faithful}'s arithmetic study,
we test \NUM{n-models} models on pairs of problems with the same equations. In one version,
a variable has an ordinary name; in the other, its name is a word for an incorrect digit.
When asked only for the final answer, models more often return the digit suggested by
the misleading name. Asking them to write the calculation before answering, known as
chain-of-thought prompting, reduces this tendency. In \NUM{eq-cot-equivalent} of
\NUM{eq-cot-cells} comparisons, the confidence interval for the added error rate lies
within \NUM{eq-delta} percentage points of zero. Writing the calculation
therefore makes answers more robust to misleading wording while the underlying math
remains unchanged.
\end{abstract}

\section{Introduction}

Consider the problem
\begin{center}
\texttt{pen=1 + cup, \quad cup=2 + 3; \quad pen=?}
\end{center}
\texttt{pen} has value 6. Renaming it \texttt{four} leaves the answer unchanged:
\begin{center}
\texttt{\textcolor{inc}{four}=1 + cup, \quad cup=2 + 3; \quad \textcolor{inc}{four}=?}
\end{center}
Only the label changes: \texttt{cup} is still 5 and the queried variable is still 6.
Answering 4 instead follows the name. We call 4 the \emph{name-suggested value} and that
answer a \emph{name error}. Because 4 occurs nowhere in the equations and differs from
every variable's value, this error can be distinguished from copying a digit.

Can asking the model to calculate before answering reduce this error? In our
\emph{answer-only} condition, the model is prompted to return just the final digit.
In our \emph{chain-of-thought} (CoT) condition, it is prompted to restate and evaluate
the equations before giving that digit. Here the calculation includes
\texttt{cup=2+3=5} and \texttt{four=1+5=6}. The comparison concerns what the model
generates; an answer-only response can still involve internal computation.

Misleading names already reduce accuracy on code tasks
\cite{le2025names,lam2025codecrash,le2026obfuscated}. We use arithmetic to make the
competing answers explicit, building on \citet{kudo2026faithful}'s study of computation
during CoT. We add a name that suggests a different answer while preserving the equations.
This conflict, like those in Stroop and shape--texture studies
\cite{stroop1935,geirhos2019texture}, separates the name's suggestion from the computed value.

A separate analysis asks which values are readable inside the model. We train a
\emph{linear probe}, a small predictor of a variable's true digit from internal activations,
on separate problems. A probe predicting 6 and the model answering 6 are separate observations.
We show that asking the model to write the calculation reduces answers matching the
misleading name in almost every tested contrast. The paired renaming comparison makes
resistance to a competing word a measurable benefit of CoT while holding the equations fixed.

\begin{figure*}[t]
\centering
\includegraphics[width=\textwidth]{figures/story_behavior.pdf}
\caption{Added answers matching the misleading name, relative to ordinary-name twins.
Bars: 95\% intervals over three demonstration sets. *Below the planned 90\% ordinary-name
accuracy floor with written calculations.}
\label{fig:master}
\end{figure*}

\section{Study design}

We use \citeauthor{kudo2026faithful}'s single-digit addition/subtraction task. The main level
requires two operations and presents the queried variable's equation first, before the
equation needed to evaluate it. Two models cover
all five levels (Appendix~\ref{app:task}). Matched versions change one name
(Table~\ref{tab:conditions}): congruent names denote the true value; incongruent names denote
another digit, absent from all inputs and variable values. We rename queried and intermediate
variables separately. All names occupy one token, preserving matched token positions.



Both prompting conditions use three fixed sets of three worked examples. We evaluate
\NUM{n-models} Llama, Gemma and OLMo checkpoints (Appendices~\ref{app:task}, \ref{app:kind}).
The main table applies the preregistered 90\% mean neutral-CoT accuracy floor; excluded
results appear in Table~\ref{tab:excluded}.

\emph{Excess name errors} subtracts how often the ordinary-name version answers the suggested
digit from how often the misleading-name version does. This controls guesses when accuracy is low.
\emph{Facilitation} is congruent-minus-neutral accuracy; \emph{interference} is
neutral-minus-incongruent accuracy. Reliable effects have consistent signs across all three
demonstration sets and 95\% intervals excluding zero. Equivalence requires a 90\% interval
inside $\pm\NUM{eq-delta}$ points. Cluster bootstraps resample matched problems with their
demonstration repeats (Appendix~\ref{app:equivalence}).

Probes train on separate neutral problems with three seeds. Random digits assigned to names
control for learning name identity \cite{hewitt2019control}. Patching replaces name activations
with a matched twin's \cite{zhang2024patching}; neutral-word and alternative-name controls
test disruption and name specificity. Most CoT measurements supply the correct calculation
token by token.
Generated-chain probes select layers on computation-disjoint validation problems before
held-out testing (Appendix~\ref{app:round6-generated}).

\section{Results}

\begin{table*}[t]
\centering\small
\resizebox{\textwidth}{!}{\tabinput{behavior}}
\caption{Level~3 effects in percentage points: means across three demonstration sets [ranges].
\emph{Helps}: congruent-name accuracy gain; \emph{Errors}: excess name errors.
$^{*}$: consistent signs and 95\% intervals excluding zero. CoT name-error contrasts meet the $\pm2$-point
equivalence bound.}
\label{tab:behavior}
\end{table*}

\subsection{Writing calculations limits misleading answers}

Figure~\ref{fig:master} shows added misleading-digit answers; Table~\ref{tab:behavior}
also reports changes in overall accuracy. Without CoT, a congruent
name raises Llama-3.1-8B's accuracy by \NUM{facilitationllama318bdirectv1} points, and a
misleading name produces \NUM{lure_excessllama318bdirectv1} points of excess name errors on the queried
variable. Across levels, the largest effects occur where direct-answer accuracy is near chance; the matched baseline prevents those arbitrary digit answers from being attributed to
the name.

With CoT, \NUM{eq-cot-equivalent} of \NUM{eq-cot-cells} name-error contrasts meet the equivalence
bound. The exception is \NUM{eq-exc-model}: its intermediate-variable excess is
\NUM{eq-exc-mean} points [\NUM{eq-exc-lo}, \NUM{eq-exc-hi}]. Its neutral CoT accuracy is
\NUM{eq-exc-neutral-acc}\%, compared with \NUM{eq-exc-letter-acc}\% under letter names.
CoT limits answers matching the misleading digit. Overall accuracy still depends on names:
\NUM{cot-acc-cells} accuracy contrasts remain reliable
(Appendix~\ref{app:equivalence}).





\subsection{Correct written values can precede name errors}
\label{sec:written-controls}

A generated calculation can state the correct value and still end with the suggested
digit. Among matched versions that both state the intermediate variable's correct value,
OLMo-2-1B still has \NUM{vw-olmo-excess} points of excess name errors
[\NUM{vw-olmo-lo}, \NUM{vw-olmo-hi}] across \NUM{vw-olmo-n} matched observations from
\NUM{vw-olmo-seeds} demonstration sets (Appendix~\ref{app:controls}). This conditional
comparison shows that stating the value does not guarantee a correct answer; it does
not estimate the causal effect of stating it.

Nor do name errors always repeat the preceding digit. Across \NUM{copy-n-err} generated
CoT name errors, pooled over models, levels and demonstration sets, only
\NUM{copy-lure-frac} have the name-suggested digit at the end of the calculation segment
before the final answer (Appendix~\ref{app:sweeps}). A name error need not copy a trailing
instance of that digit.

We also test writing values alone, such as \texttt{cup=5, four=6}, instead of restating
and evaluating equations. On two-operation problems, ordinary-name accuracy falls to
\NUM{simple-l8b-acc}\% for Llama-3.1-8B and \NUM{simple-l3b-acc}\% for Llama-3.2-3B
(Appendix~\ref{app:simple}). This format changes both arithmetic performance and the
written context. It cannot isolate an effect of name repetition while preserving computation.

\subsection{What the probes recover}

These probes read internal states while the correct calculation is supplied. Just after
the intermediate equation, Llama-3.2-3B's correct-digit accuracy is \NUM{p2-acc-pct}\%,
\NUM{p2-sel-points} points above the name-identity control. At the query, the control reaches
\NUM{p3-ctl-pct}\%, showing that the probe can identify the name. Before the value is
written, correct-digit accuracy reaches \NUM{p4-acc-pct}\% and predictions of the
name-suggested digit fall to zero. Gemma's two variables show different readout timing
(Appendix~\ref{app:probe-extensions}).

Across \NUM{competence-models} models, \NUM{vs-decodable} of \NUM{vs-cells} variable readouts
have neutral accuracy of \NUM{vs-decodable-min-pct}--\NUM{vs-decodable-max-pct}\% and excess
name errors within \NUM{vs-decodable-maxexcess} points of zero. OLMo-2-1B's intermediate
readout scores \NUM{vs-undecodable-acc-pct}\%. But the correlation between each model's mean
misleading-name probe accuracy and excess name errors reverses from
$r=\NUM{competence-probe-lure-r}$ to
$\NUM{competence-probe-lure-without-olmo-r}$ when OLMo-2-1B is omitted
(Appendix~\ref{app:round5-competence}). Overall task competence remains a competing account.
Within-problem fits also disagree on the association's sign (Appendix~\ref{app:sweeps}).

Probes trained on correct calculations score \NUM{generated-olmo1-cal-lo}--\NUM{generated-olmo1-cal-hi}\% on
OLMo-2-1B's neutral generated chains. Training on generated neutral chains raises this to
\NUM{r6-generated-olmo1-cal-min}--\NUM{r6-generated-olmo1-cal-max}\%, still below the 90\%
minimum for interpreting these probes in every cell (Appendix Figure~\ref{fig:readout-check}). Llama-3.2-3B reaches
\NUM{r6-generated-llama3-cal-min}--\NUM{r6-generated-llama3-cal-max}\% with the same design,
but its name-error cases remain rare. OLMo's generated-chain readouts therefore
remain descriptive (Appendix~\ref{app:round6-generated}).

\subsection{Patching localizes the answer-only effect}

For Llama-3.2-3B without CoT, replacing name activations with the neutral twin's removes
\NUM{lure-removed-3b} of name errors when all layers are patched (Figure~\ref{fig:patch}).
A span-by-layer sweep places the strongest effects at the
defining equation in layers \NUM{grid-layers}, where removal reaches \NUM{grid-top-removed}
(Appendix~\ref{app:figs}). The neutral-word control also damages \NUM{control-damage} of correct
answers. The patch identifies a region sensitive to the name; disruption limits claims about the restored computation.

The Llama-3.1-8B controls do not distinguish a name-specific cause: neutral and alternative-name
patches remove errors similarly. Under a gold CoT, injecting the misleading name into a neutral
run of the exception model moves \NUM{exc-inject} of answers to the name-suggested value, while the word control
changes \NUM{exc-ctlword} at the same layer. These interventions test only the selected models and positions; they do not establish a common causal mechanism across models.



\section{Implications and next experiments}

Writing calculations improves resistance to a word that suggests the wrong answer while
the mathematical problem stays fixed. This gives reasoning evaluations a concrete test
beyond aggregate accuracy: rename a variable and measure whether answers move toward
the implied digit. The remaining accuracy differences show why both outcomes matter.
The written-value control also distinguishes a correct value in the calculation from
a correct final answer: seeing that value alone can miss the effect of a misleading name.

Two experiments would help explain the reduction. First, vary whether the response
repeats variable names and whether it includes equations, using formats matched for
length and ordinary-name accuracy. The values-only result motivates this matching:
otherwise a change in name errors may accompany a loss of arithmetic competence.
Second, hold a correct generated calculation fixed, vary whether the final query repeats
the misleading name, and request either a digit alone or a named assignment. This would
test whether the misleading cue can redirect the answer after the correct value has
been written. Both are proposed
tests; the present comparisons do not identify which part of CoT produces the benefit.

\section{Related work}

\citet{kudo2026faithful} locate intermediate-value readouts during arithmetic CoT; our competing
names distinguish a computed value from a lexical cue. Recent language-model Stroop studies
test lexical responses \cite{wang2026stroop,hu2026conflict}. Faithfulness studies perturb prompts, chains or
parameters \cite{turpin2023unfaithful,lanham2023measuring,chen2025reasoning,tutek2025unlearning};
we combine a known conflicting cue with internal readouts. Positional copying and output priors
can also explain arithmetic errors \cite{liu2026readout,prabhakar2024deciphering};
Section~\ref{sec:written-controls} checks repetition of the digit preceding the answer.
A companion code study \citep{anon2026code} tests readouts before actual wrong value writes
in Python traces and changes the demonstration format. Here the main finding is a reduction
in arithmetic answers matching a misleading name while the equations remain unchanged.

\section{Conclusion}

Writing and evaluating equations limits answers suggested by misleading names in almost
every tested arithmetic contrast. Paired renaming measures this benefit directly;
the written-value and probe controls guide further tests of how it arises.

\section*{Limitations}

\begin{itemize}
\item \textbf{Scope.} Short synthetic arithmetic with English number words does not establish
real-code effects. Identifier-style names cover one level. Models are at most 12B, with raw
completion prompts; chat templates and long-form reasoning are untested. Accuracy varies
with demonstrations, so behavioral claims require agreement across three sets.
\item \textbf{Internal measurements.} Probes establish decodability, not use. Most measurements
supply correct chains. Generated-chain calibration passes for Llama-3.2-3B but fails for
OLMo-2-1B, even after retraining (Appendix~\ref{app:round6-generated}); only two models are tested.
The nine-model correlation is descriptive, reverses without OLMo, and confounds value
availability with competence; variable roles are not independent models.
\item \textbf{Interventions.} Patches replace mixed representations and damage correct answers.
CoT also repeats names and equations; the value-only format changes performance and is not
an equivalent intervention (Appendix~\ref{app:simple}).
\end{itemize}

\bibliography{refs}

\appendix

\section{Task, models and replication}
\label{app:task}

\begin{table}[t]
\centering\small
\resizebox{\columnwidth}{!}{\begin{tabular}{@{}ll@{}}
\toprule
Name & Problem (answer: 6) \\
\midrule
Neutral & \texttt{pen=1 + cup, cup=2 + 3; pen=?} \\
Congruent & \texttt{\textcolor{con}{six}=1 + cup, cup=2 + 3; six=?} \\
Incongruent & \texttt{\textcolor{inc}{four}=1 + cup, cup=2 + 3; four=?} \\
\bottomrule
\end{tabular}}
\caption{Matched names preserve the arithmetic. The misleading digit, 4, occurs nowhere else.}
\label{tab:conditions}
\end{table}
A problem contains single-digit variables, addition or subtraction, and a query. The five
levels vary dependency order, depth and the presence of an unused equation
(Table~\ref{tab:levels}). Three same-level worked examples precede each test problem.
The evaluated models are \NUM{model-list}, plus \NUM{n-ladder} derived checkpoints
(Appendix~\ref{app:kind}).

\begin{table}[t]
\centering\small\setlength{\tabcolsep}{4pt}
\resizebox{\columnwidth}{!}{\begin{tabular}{@{}clccc@{}}
\toprule
Level & Problem (neutral names) & Steps & Held & Distr. \\
\midrule
1 & \texttt{pen=1 + cup, cup=2; pen=?} & 1 & 1 & 0 \\
2 & \texttt{pen=2 + 3, cup=1 + pen; cup=?} & 2 & 0 & 0 \\
3 & \texttt{pen=1 + cup, cup=2 + 3; pen=?} & 2 & 1 & 0 \\
4 & \texttt{pen=1 + cup, cup=2 + 3,} & 2 & 1 & 1 \\
  & \texttt{\ \ box=4 + 5; pen=?} & & & \\
5 & \texttt{pen=1 + cup, cup=2 + box,} & 3 & 2 & 0 \\
  & \texttt{\ \ box=1 + 2; pen=?} & & & \\
\bottomrule
\end{tabular}}
\caption{The five levels of \citet{kudo2026faithful}: additions or subtractions needed
(\emph{Steps}), variables unresolved when first read (\emph{Held}), and unused equations
(\emph{Distr.}). At level~1 the intermediate value is stated rather than computed.}
\label{tab:levels}
\end{table}

We reimplement \citet{kudo2026faithful}'s task and probe recipe rather than using their code.
The shared choices are single-digit outputs, same-level demonstrations, a 10-way linear
residual-stream classifier, full-batch SGD at $10^{-3}$ for 10,000 epochs, and their
span-by-layer-window patching design. For replication we train on letter names and report
$t^{*}$, $t^{*}_{\mathrm{eq}}$, $\mathrm{Acc}^{\prec\mathrm{CoT}}$ and
$\mathrm{Acc}^{\succ\mathrm{CoT}}$ at $\tau=0.9$.
The Llama-3.2-3B value-writing positions agree with their Table~3: equation
\NUM{replication-eq-v1} for the queried variable and \NUM{replication-eq-v2} for the intermediate.
Our additions are the conflicting names, matched baselines, demonstration-set claim rule and
equivalence bound.



On the letter condition, pre-chain probe maxima for Llama-3.2-3B are
\NUM{replication-pre-v1} and \NUM{replication-pre-v2}, compared with 17.8 and 33.2 reported by
\citet{kudo2026faithful}. The value-writing positions agree; the accuracy difference may reflect
training details. Llama-3.1-8B's intermediate value becomes decodable at equation
\NUM{replication8b-eq-v2}, but its queried value becomes decodable at equation
\NUM{replication8b-eq-v1}, earlier than Llama-3.2-3B's equation \NUM{replication-eq-v1}.
The larger model can therefore expose that value before writing its final evaluation.
Full heatmaps appear in Appendix~\ref{app:figs}.

\begin{table*}[t]
\centering\small
\tabinput{behavior_excluded}
\caption{The level-3 model below the preregistered 90\% mean neutral-CoT accuracy floor,
excluded from the main behavioral table. Columns and claim markers follow Table~\ref{tab:behavior}.}
\label{tab:excluded}
\end{table*}

\section{Equivalence checks}
\label{app:equivalence}
Following preregistration Amendment~3, equivalence requires a 90\% cluster-bootstrap
interval wholly inside the $\pm\NUM{eq-delta}$-point margin (two one-sided 5\% tests).
We resample matched sets with their demonstration-seed replicates together, using 4,000
bootstrap draws. Other behavioral intervals remain at 95\%.
For CoT accuracy, \NUM{eq-acc-equivalent} of \NUM{eq-acc-cells} facilitation and interference
contrasts meet the same equivalence bound. Of the \NUM{cot-acc-cells} reliable accuracy
contrasts, \NUM{eq-acc-claimable-equivalent} are also equivalent within the margin and
\NUM{eq-acc-claimable-outside} are not; a detectable effect can still be small.
The written-value control uses the same matched pairs for its excess estimate and interval;
the exception model's estimate pools all three demonstration sets and clusters on matched set.
\begin{table}[h]
\centering\small
\tabinput{equivalence}
\caption{Equivalence within the two-point margin uses 90\% intervals; reliability uses
95\% intervals excluding zero and the same sign across all three demonstration sets.
Each accuracy cell is one facilitation or interference contrast.}
\end{table}
\IfFileExists{tables/exception_patching.tex}{%
\begin{table*}[t]
\centering\small
\resizebox{\textwidth}{!}{\tabinput{exception_patching}}
\caption{OLMo-2-1B-Instruct, CoT, \NUM{exc-controls-n} matched pairs per target, patching name tokens at all
layers. Removed: percentage of baseline name-suggested predictions removed by the neutral-twin patch.
Word damage: percentage of correct neutral predictions damaged by the word control.
Alt. removed: baseline name-suggested predictions removed by the alternative-name patch. Alt. followed:
percentage following that source's name-suggested value. Rates are percentages;
-- means undefined. These readouts force the gold chain, so their baseline errors differ
from errors in the model's own generated chains.}
\end{table*}
}{}
\section{Tokenizer checks}
\label{app:tok}
We test candidate names in five tokenizer contexts and retain words that form one token
without merging with a neighbor. For Llama~3, all ten number words and 95 of 112 neutral
candidates pass. Gemma and OLMo use the same lists; runtime checks confirm equal token lengths
and positions within matched groups. Spaces around binary operators prevent merges such as
\texttt{1-two}. No instance in the reported runs is excluded by the layout check.

\section{Additional naming and written-value controls}
\label{app:controls}

\paragraph{Written values.} We split generated chains by whether they state the target's true
value. An observation enters the excess estimate only when its neutral twin falls in the same
split. The estimate and its 95\% interval use those same paired observations. Pooled over
eligible cells, the excess is \NUM{vw-excess-wrote} points in \NUM{vw-cells-wrote} cells where
the value is stated and \NUM{vw-excess-not} in \NUM{vw-cells-not} where it is not. The exception
persists in the stated-value split (main text), so this split does not explain its effect.

\paragraph{Identifier-style names.} At level~3, \NUM{ident-models} models receive names such as
\texttt{q4} instead of \texttt{four}, under three demonstration sets. Without CoT, queried-variable
excess name errors range from \NUM{ident-direct-lure-lo} to \NUM{ident-direct-lure-hi} points;
\NUM{ident-direct-lure-claim} of \NUM{ident-direct-cells} cells are reliable. For number words
on the same models, the range is \NUM{word-direct-lure-lo} to \NUM{word-direct-lure-hi}, with
\NUM{word-direct-lure-claim} of \NUM{word-direct-cells} reliable. The numeric name-error effect is
therefore stronger for number words in this comparison; congruent identifier names can still
improve accuracy.

\paragraph{Word-class effects.} With CoT at level~4, a number word on the queried variable costs
Llama-3.2-3B \NUM{wordclass-cost-L4} accuracy points whether the name agrees with the value or
not. Those errors restate the wrong equation. This is distinct from answering the digit the
name denotes, which is why the name-suggested value result cannot be read as removal of every naming effect.

\paragraph{Unused variables.} A number word on the level~4 distractor changes accuracy by at
most \NUM{irr-acc-delta-max} points. Its no-CoT excess name errors range from \NUM{irr-direct-lo} to
\NUM{irr-direct-hi}; \NUM{irr-direct-claimable} of \NUM{irr-direct-cells} cells are reliable.
With CoT, the range is \NUM{irr-cot-lo} to \NUM{irr-cot-hi}, with
\NUM{irr-cot-claimable} of \NUM{irr-cot-cells} reliable. At the answer position, probe probability of the name-suggested value
is \NUM{bound-mass} for a computed variable and \NUM{unbound-mass} for the distractor; the
latter's true value has probe accuracy \NUM{distractor-acc}.

\section{Full sweeps and selectivity}
\label{app:sweeps}

\begin{figure*}[t]
\centering
\IfFileExists{figures/fig1_master_L3.pdf}{\includegraphics[width=\textwidth]{figures/fig1_master_L3.pdf}}{\fbox{\parbox{0.95\textwidth}{\centering\small\textcolor{red}{[Figure 1 not generated]}}}}
\caption{Level~3 queried-variable effects: \textbf{(a)} congruent-name accuracy gains;
\textbf{(b)} excess name errors. Means and 95\% intervals across three demonstration sets;
grey: $\pm2$ points, a reference band rather than an equivalence test. \textbf{(c)} Best-layer true-value accuracy and name-suggested predictions
under a supplied correct calculation. OLMo-2-1B is below the planned neutral-accuracy floor. Shading: name positions; dotted line: position before the queried value.}
\label{fig:detailed-master}
\end{figure*}

\begin{figure}[t]
\centering
\IfFileExists{figures/fig3_patching_llama32-3b_L3_v1.pdf}{\includegraphics[width=\columnwidth]{figures/fig3_patching_llama32-3b_L3_v1.pdf}}{\fbox{\textcolor{red}{[Patching figure not generated]}}}
\caption{Llama-3.2-3B, answer-only: queried-variable name errors removed by layer, with
neutral-word damage and alternative-name controls. Removal counts changes away from the
name-suggested digit.}
\label{fig:patch}
\end{figure}

\begin{figure*}[h]
\centering
\IfFileExists{figures/fig2_tokens_llama32-3b_L3_cot_v2.pdf}{\includegraphics[width=\textwidth]{figures/fig2_tokens_llama32-3b_L3_cot_v2.pdf}}{\fbox{\parbox{0.95\textwidth}{\centering\small\textcolor{red}{[Figure 2, generated by scripts/58\_token\_figure.py]}}}}
\caption{Llama-3.2-3B, intermediate variable, misleading-name problems. Top: best-layer
correct-digit and name-suggested predictions; bottom: correct-digit accuracy by layer.
The correct calculation is supplied. Rates pool problems; tokens illustrate one example.}
\label{fig:heat}
\end{figure*}

\paragraph{Mixed-effects check.} We compare the paired bootstrap with a mixed logit model,
\texttt{name\_error \textasciitilde{} incongruent + C(seed)}, with a random intercept for
matched set. Each twin's outcome refers to the same name-suggested digit. Both procedures use the same
fixed subsample of \NUM{mix-cap} sets per cell. Of \NUM{mix-cells} fitted cells,
\NUM{mix-agree} agree on whether the interval excludes zero; \NUM{mix-unfitted} additional
cells have no outcome variation. The mixed fit supports an effect in \NUM{mix-gained} cells
where the bootstrap does not and withdraws support in \NUM{mix-lost} cells where it does
(\NUM{mix-lost-cells}). Subsampling also removes support from \NUM{mix-sub-lost-n} starred
Table~\ref{tab:behavior} cell under both procedures (\NUM{mix-sub-lost-cells}). The mixed
model uses variational Bayes, so its interval is an approximate posterior interval. We retain
the full-data bootstrap as the primary estimate on the reported percentage-point scale.

\paragraph{Instance-level association.} Within each model, level, regime, variable and
position, we regress name errors on the probe margin, $\log p_{\mathrm{true}}-\log p_{\mathrm{name}}$,
where $p_{\mathrm{name}}$ is the probe probability of the name-suggested value,
with standard errors clustered on matched set. The layer is selected by neutral accuracy.
Of \NUM{link-fits} fits, \NUM{link-sig} have $p<.05$: \NUM{link-neg} coefficients are negative
and \NUM{link-pos} positive. The direction is not consistent, and \NUM{link-pos-topmodel-n}
positive results come from the model supplying most fits. This leaves the instance-level
association unresolved. The model-level pattern is also sensitive to panel composition
(Appendix~\ref{app:round5-competence}).

\paragraph{Name-error share.} For comparison with shape--texture studies \cite{geirhos2019texture},
we also divide name errors by the total of name errors and correct answers, excluding congruent trials.
Across \NUM{lureindex-direct-n} level~3 model--variable pairs, the index ranges from
\NUM{lureindex-direct-lo} to \NUM{lureindex-direct-hi} without CoT and from
\NUM{lureindex-cot-lo} to \NUM{lureindex-cot-hi} with CoT. The main text uses paired excess
because the index does not subtract how often the neutral model answers that digit.

\paragraph{Positional-copy control.} Across \NUM{copy-n-err} generated CoT name errors over
models, levels and demonstration sets, the number immediately preceding the answer is the
name-suggested value in \NUM{copy-lure-frac} of cases and the target's true value in \NUM{copy-true-frac}.
A trailing name-suggested value is therefore absent from many errors. This check constrains the positional-copy
explanation of \citet{liu2026readout}, but does not identify what was copied in each case.

Probes report true-value accuracy, name-suggested prediction rate, probability of the name-suggested value and margin, with name-identity
control accuracy and selectivity \cite{hewitt2019control}. P1 is the name token, P2 the end of
its defining equation, P3 the query, P4 the position before its value is written, and P5 the
position before the final answer. P1 is a reference, not evidence for a computed value.
At P3 the control is near~1, so a readout of the name-suggested value there may reflect name identity.
Figure~\ref{fig:heat} shows the CoT token sweep; Appendix~\ref{app:figs} contains the other
conditions and direct regime. Letter-trained replication uses the metrics in
Appendix~\ref{app:task}.

\begin{figure}[h]
\centering
\IfFileExists{figures/fig5_kudo3_llama32-3b_L3_cot_v1_lureerr_i0.pdf}{\includegraphics[width=\columnwidth]{figures/fig5_kudo3_llama32-3b_L3_cot_v1_lureerr_i0.pdf}}{\fbox{\parbox{0.95\columnwidth}{\centering\small\textcolor{red}{[Figure 4, generated by scripts/64\_kudo\_fig3.py]}}}}
\caption{Correct-calculation probe predictions for a problem answered with the
name-suggested digit. Top: neutral modal prediction (one seed); bottom: majority prediction
(three seeds). Green: true value; red: wrong answer; amber: name-suggested value; grey: other.}
\label{fig:kudo3}
\end{figure}

\begin{figure*}[h]
\centering
\IfFileExists{figures/fig_own_chain_llama32-3b_L3_cot_v1.pdf}{\includegraphics[width=\textwidth]{figures/fig_own_chain_llama32-3b_L3_cot_v1.pdf}}{\fbox{\parbox{0.95\textwidth}{\centering\small\textcolor{red}{[own-chain figure, generated by scripts/67\_own\_chain\_figure.py]}}}}
\caption{Value-step probes applied to Llama-3.2-3B's generated calculations on the first
\NUM{ownchain-n} of \NUM{ownchain-total} incorrectly answered problems. Each panel is one
problem; colours follow the legend. Training uses supplied correct calculations on
ordinary-name problems.}
\label{fig:ownchain}
\end{figure*}

\section{How the models were tuned}
\label{app:kind}

We compare six checkpoints on the Llama-3.1-8B architecture: pretrained, instruction-tuned,
single-task LoRA, multi-task LoRA, TIES-merged adapters and reasoning-tuned. Adapters are folded
into the weights before evaluation. Shared tokenizers and identical matched problems permit
paired comparisons (Table~\ref{tab:kind}). A pretrained/instruct Gemma pair supplies a second
family comparison.

\begin{table}[h]
\centering\small
\resizebox{\columnwidth}{!}{\tabinput{modelkind}}
\caption{Level 3, no CoT, name errors above the matched twin's rate, mean over three
demonstration sets. \emph{Acc.} is neutral accuracy. $^{*}$ as in Table~\ref{tab:behavior}.
The bottom block repeats the pretrained-to-instruct step in a second family, Gemma-3-4B.}
\label{tab:kind}
\end{table}

The pretrained and reasoning-tuned Llama checkpoints have reliable positive queried-variable
excess name errors without CoT (\NUM{kind-base-v1} and \NUM{kind-reas-v1} points). The three task-tuned
variants differ from the instruction-tuned checkpoint by at most \NUM{kind-maxgap} points
across \NUM{kind-tuned-cells} comparisons. These are observed differences, not tests of
checkpoint-to-checkpoint equivalence. Every checkpoint's CoT excess is at most
\NUM{kind-cot-max} points in magnitude.

Task accuracy complicates interpretation. The multi-task model's reliable negative excess
(\NUM{kind-ftm-v1}) accompanies lower neutral accuracy (\NUM{kind-ftm-acc}\%, versus
\NUM{kind-it-acc}\%). The reasoning-tuned checkpoint has the lowest accuracy
(\NUM{kind-reas-acc}\%). It also requires a system instruction to enable long-form reasoning;
our raw-completion prompts leave that mode off. The result concerns the checkpoint under this
prompt, not reasoning mode in deployment.

Only the instruct Gemma effect passes the claim rule (\NUM{kind-gi-v1}, versus
\NUM{kind-gb-v1} pretrained). The difference is untested and both estimates lie within the
margin, so the pair cannot establish where in training name interference arises.

\section{Probe recipe and patching scope}
\label{app:recipe}

We check two implementation choices inherited from \citet{kudo2026faithful}: the probe optimizer and the name occurrences included in a patch.

\paragraph{Recipe.} We refit the level-3 probes for Llama-3.2-3B with L-BFGS, and with SGD on
standardized inputs, against the reported recipe.

\begin{table}[h]
\centering\small
\resizebox{\columnwidth}{!}{\tabinput{recipe}}
\caption{Best-layer true-value accuracy on neutral problems under three probe recipes
(proportions from 0 to 1).}
\label{tab:recipe}
\end{table}

All three recipes give the same qualitative result: the value is decodable where the chain writes it
(\NUM{e17-value-lo}--\NUM{e17-value-hi}) and not at the end of its defining equation
(\NUM{e17-end-lo}--\NUM{e17-end-hi}). The query position moves, from \NUM{e17-sgd-query-v2}
under the reported recipe to \NUM{e17-std-query-v2} under standardization, within a range of
\NUM{e17-query-lo}--\NUM{e17-query-hi} over recipes and variables. We already report that
position as the presence of the name rather than a decoded value, because its control task
scores \NUM{p3-ctl}, so this variation does not change the interpretation of that position.

\paragraph{Scope.} Patching every occurrence of the name includes the copy the model writes in
its own output. In the direct regime, restricting the patch to the name's occurrences in the prompt leaves the
intermediate variable essentially unchanged, since its name never appears in the answer: removal is
\NUM{e18-v2-all-removed}\% at layers~\NUM{e18-v2-all-best} under both scopes. For the queried
variable the best windows remove similar shares (\NUM{e18-v1-prompt-removed}\% at
layers~\NUM{e18-v1-prompt-best} prompt-only, against \NUM{e18-v1-all-removed}\% at
layers~\NUM{e18-v1-all-best}), the prompt-only patch removes less in the later windows, and the
two disagree most in the first window, where the prompt-only patch removes
\NUM{e18-v1-prompt-w03}\% and damages \NUM{e18-v1-prompt-w03damage}\% of correct answers. That
is mechanical: the queried variable's name also sits at the position we read, so patching its
prompt copies alone leaves the model reading a name whose early representation no longer
matches. The localization we report does not depend on the choice.

\section{A chain that writes values without equations}
\label{app:simple}

Our chain restates each equation before substituting into it, as
\citeauthor{kudo2026faithful}'s does. Their Tables~5--6 also report a format that writes only the
sub-results, \texttt{cup=5, pen=6}. That format is the arithmetic counterpart of a code trace,
which writes a variable's value and never the expression that produced it, so we ran it on the
panel at level 3 with three demonstration sets, with two predictions registered in advance for both Llama models: that
the excess name errors at the intermediate's value step would exceed the \NUM{eq-delta}-point
margin, and that the value would still be decodable there (probe accuracy at least 0.85), making
any failure one of readout.

Both predictions fail as registered. Llama-3.1-8B writes the name-suggested value at the intermediate's value
step \NUM{simple-l8b-v2} points above its twin, inside the margin. Llama-3.2-3B is at
\NUM{simple-l3b-v2}, which passes our claim rule, but across the three demonstration sets that
is \NUM{simple-l3b-v2-seeds}, so one set carries most of it. More informative is what the probes say. At
the step that writes the value, the true value is decodable at \NUM{simpleprobe-l8b-simple-v1}
and \NUM{simpleprobe-l8b-simple-v2} for Llama-3.1-8B and at
\NUM{simpleprobe-l3b-simple-v1} and \NUM{simpleprobe-l3b-simple-v2} for Llama-3.2-3B, against
\NUM{simpleprobe-l8b-cot-v1} and \NUM{simpleprobe-l8b-cot-v2} for the same model under the full
chain. Neutral accuracy falls with it, to \NUM{simple-l8b-acc}\% and \NUM{simple-l3b-acc}\%.

Removing the equations lowers both accuracy and true-value probe readouts. It therefore fails to isolate a readout error while preserving those measures. This result does not support a shared explanation for value-only arithmetic chains and code traces.

\section{Additional levels, models and demonstration sets}
Table~\ref{tab:levels-appendix} repeats Table~\ref{tab:behavior} for the other four levels.
The pretrained Gemma-3-4B, the base half of the Gemma pair, is reported in
Table~\ref{tab:kind} rather than in Table~\ref{tab:behavior}.

\paragraph{Data, compute and software.} Each level has \NUM{n-test-sets} matched sets, each
rendered in every condition, and \NUM{n-probe-train} separately generated problems for
training probes; every SGD probe is trained with \NUM{probe-seeds} seeds. Every run here is
inference and the probes are linear; the three adapter checkpoints of Appendix~\ref{app:kind}
were trained in a separate project, and that compute is not counted below. The recorded allocation time totals at most \NUM{compute-gpuh} GPU-hours on NVIDIA A30, H100 and H200 cards, including idle time and shared workers that also ran companion-code experiments.
The models are from the Llama~3 \citep{grattafiori2024llama3}, Gemma~3 \citep{gemmateam2025gemma3}
and OLMo~2 \citep{olmo2024olmo2} families, with the reasoning-tuned checkpoint from
\citet{bercovich2025nemotron}, each used under its release licence for research. Models run in
PyTorch \citep{paszke2019pytorch} with Transformers \citep{wolf2020transformers}, and the probes are
trained in PyTorch; the L-BFGS refit of Appendix~\ref{app:recipe} uses scikit-learn
\citep{pedregosa2011sklearn}, and the mixed-effects check uses statsmodels
\citep{seabold2010statsmodels}.

\begin{table}[h]
\centering\small
\resizebox{\columnwidth}{!}{\tabinput{behavior_levels}}
\caption{As Table~\ref{tab:behavior}, for levels 1, 2, 4 and 5 (Llama-3.2-3B and Llama-3.1-8B). Level~1 is the no-computation
control: the intermediate variable's value is stated rather than computed.}
\label{tab:levels-appendix}
\end{table}

Per-set values, the pooled estimate and the claim rule are released with the code, together
with the level-4 irrelevant-name-suggested value
control, the name-suggested value-distance analysis, the teacher-forced name-suggested prediction rates, and the error taxonomy
behind the word-class effect.

% Round-four probe extensions
\section{Additional value readouts}
\label{app:probe-extensions}

We extend the value-step comparison to Llama-3.2-3B-Instruct, Llama-3.1-8B-Instruct and
Gemma-3-12B-Instruct. Each probe trains on 10,000 separate neutral problems with three probe
seeds, using demonstration set 7. Table~\ref{tab:probe-extensions} reports the added pairs;
the main comparison includes them. We select layers by neutral selectivity, as in the original
analysis. The misleading-name accuracy uses the same selected layer. This extension tests
the aggregate association; it does not resolve the within-problem link.

\begin{table*}[t]
\centering\small
\tabinput{probe_extensions}
\caption{Additional probes at each variable's value-writing position. Neutral and misleading
columns are true-digit accuracy; control is neutral name-identity accuracy. Rates are percentages,
averaged over three probe seeds. The layer maximizes neutral accuracy minus control accuracy.}
\label{tab:probe-extensions}
\end{table*}

We also probe every token for Gemma-3-4B-Instruct (Figures~\ref{fig:gemma-query}
and~\ref{fig:gemma-intermediate}), using 4,000 independent neutral training problems and
1,000 test problems per condition. The tokens shown come from one example; the curves and
heatmaps pool problems with different names and values. Maximum-over-layers curves describe
where a digit can be read; they do not identify the representation used to answer.
On neutral problems, the best pre-chain accuracy is \NUM{gemma-token-pre-v1}\% for the
queried digit and \NUM{gemma-token-pre-v2}\% for the intermediate digit. The queried readout
\NUM{gemma-token-timing-v1}; the intermediate readout \NUM{gemma-token-timing-v2}.
Token zero is the first token of the chain; negative positions are in the prompt.

\begin{figure*}[t]
\centering
\includegraphics[width=\textwidth]{figures/fig2_tokens_gemma3-4b-it_L3_cot_v1.pdf}
\caption{Gemma-3-4B-Instruct, queried variable, misleading-name problems under the correct
chain. Top: maximum-over-layers true-value accuracy and name-suggested prediction rate; bottom: true-value accuracy
by layer. Shading marks the variable's name; the dotted line marks the position before its
value; the arrow marks the supplied digit itself.}
\label{fig:gemma-query}
\end{figure*}

\begin{figure*}[t]
\centering
\includegraphics[width=\textwidth]{figures/fig2_tokens_gemma3-4b-it_L3_cot_v2.pdf}
\caption{The same Gemma token sweep for the intermediate variable. Both figures use
neutral-trained probes; no misleading-name problems enter training.}
\label{fig:gemma-intermediate}
\end{figure*}

% generated by scripts/66_figure_dump.py
\clearpage
\section{Every figure at full size}
\label{app:figs}
Each figure below is reproduced as a full-width figure on a portrait page, within the ACL margins.

\begin{figure*}[p]\centering
\includegraphics[width=\textwidth,height=0.85\textheight,keepaspectratio]{figures/fig1_master_L3.pdf}
\caption{Queried-variable effects at level 3: congruent-name accuracy gains, excess name errors, and best-layer correct-calculation readouts. Bars: 95\% intervals; grey: a two-point reference band. OLMo-2-1B is below the planned accuracy floor.}\label{fig:dump-fig1-master-L3}
\end{figure*}
\clearpage

\begin{figure*}[p]\centering
\includegraphics[width=\textwidth,height=0.85\textheight,keepaspectratio]{figures/fig2_tokens_llama32-3b_L3_cot_v2.pdf}
\caption{Per-token probes, written-calculation regime, intermediate variable.}\label{fig:dump-fig2-tokens-llama32-3b-L3-cot-v2}
\end{figure*}
\clearpage

\begin{figure*}[p]\centering
\includegraphics[width=\textwidth,height=0.85\textheight,keepaspectratio]{figures/fig2_tokens_llama32-3b_L3_cot_v1.pdf}
\caption{Per-token probes, chain-of-thought regime, queried variable.}\label{fig:dump-fig2-tokens-llama32-3b-L3-cot-v1}
\end{figure*}
\clearpage

\begin{figure*}[p]\centering
\includegraphics[width=\textwidth,height=0.85\textheight,keepaspectratio]{figures/fig2_tokens_llama32-3b_L3_direct_v2.pdf}
\caption{Per-token probes, direct regime, intermediate variable.}\label{fig:dump-fig2-tokens-llama32-3b-L3-direct-v2}
\end{figure*}
\clearpage

\begin{figure*}[p]\centering
\includegraphics[width=\textwidth,height=0.85\textheight,keepaspectratio]{figures/fig4_errors_llama32-3b_L3_cot_v1.pdf}
\caption{Probing accuracy split by whether the model answers the problem correctly on its own, with the layer sweep for each half.}\label{fig:dump-fig4-errors-llama32-3b-L3-cot-v1}
\end{figure*}
\clearpage

\begin{figure*}[p]\centering
\includegraphics[width=\textwidth,height=0.85\textheight,keepaspectratio]{figures/fig5_kudo3_llama32-3b_L3_cot_v1_lureerr_i0.pdf}
\caption{Top-1 probe prediction at every layer and token, on a problem answered with the name-suggested value. Layout after \citeauthor{kudo2026faithful}'s Figure~3.}\label{fig:dump-fig5-kudo3-llama32-3b-L3-cot-v1-lureerr-i0}
\end{figure*}
\clearpage

\begin{figure*}[p]\centering
\includegraphics[width=\textwidth,height=0.85\textheight,keepaspectratio]{figures/fig5_kudo3_llama32-3b_L3_cot_v1_lureerr_i1.pdf}
\caption{Queried-variable example answered with the digit suggested by its misleading name.}\label{fig:dump-fig5-kudo3-llama32-3b-L3-cot-v1-lureerr-i1}
\end{figure*}
\clearpage

\begin{figure*}[p]\centering
\includegraphics[width=\textwidth,height=0.85\textheight,keepaspectratio]{figures/fig5_kudo3_llama32-3b_L3_cot_v1_lureerr_i2.pdf}
\caption{Queried-variable example answered with the digit suggested by its misleading name.}\label{fig:dump-fig5-kudo3-llama32-3b-L3-cot-v1-lureerr-i2}
\end{figure*}
\clearpage

\begin{figure*}[p]\centering
\includegraphics[width=\textwidth,height=0.85\textheight,keepaspectratio]{figures/fig5_kudo3_llama32-3b_L3_cot_v1_err_i0.pdf}
\caption{Queried-variable example whose generated final answer is wrong and differs from the name-suggested digit.}\label{fig:dump-fig5-kudo3-llama32-3b-L3-cot-v1-err-i0}
\end{figure*}
\clearpage

\begin{figure*}[p]\centering
\includegraphics[width=\textwidth,height=0.85\textheight,keepaspectratio]{figures/fig5_kudo3_llama32-3b_L3_cot_v1_err_i1.pdf}
\caption{Queried-variable example whose generated final answer is wrong and differs from the name-suggested digit.}\label{fig:dump-fig5-kudo3-llama32-3b-L3-cot-v1-err-i1}
\end{figure*}
\clearpage

\begin{figure*}[p]\centering
\includegraphics[width=\textwidth,height=0.85\textheight,keepaspectratio]{figures/fig5_kudo3_llama32-3b_L3_cot_v2_err_i0.pdf}
\caption{Intermediate-variable example from a problem answered wrongly. Red is the model's final answer, which is the queried variable's value, so a cell that matches it is not a readout of this variable.}\label{fig:dump-fig5-kudo3-llama32-3b-L3-cot-v2-err-i0}
\end{figure*}
\clearpage

\begin{figure*}[p]\centering
\includegraphics[width=\textwidth,height=0.85\textheight,keepaspectratio]{figures/fig5_kudo3_llama32-3b_L3_cot_v2_err_i1.pdf}
\caption{Intermediate-variable example from a problem answered wrongly.}\label{fig:dump-fig5-kudo3-llama32-3b-L3-cot-v2-err-i1}
\end{figure*}
\clearpage

\begin{figure*}[p]\centering
\includegraphics[width=\textwidth,height=0.85\textheight,keepaspectratio]{figures/kudo_fig2_llama32-3b_L3_cot_v2.pdf}
\caption{Accuracy heatmaps by token and layer for the neutral, congruent and incongruent conditions, in the layout of \citeauthor{kudo2026faithful}'s Figure~2.}\label{fig:dump-kudo-fig2-llama32-3b-L3-cot-v2}
\end{figure*}
\clearpage

\begin{figure*}[p]\centering
\includegraphics[width=\textwidth,height=0.85\textheight,keepaspectratio]{figures/kudo_fig2_llama32-3b_L3_direct_v2.pdf}
\caption{The same heatmaps in the direct regime.}\label{fig:dump-kudo-fig2-llama32-3b-L3-direct-v2}
\end{figure*}
\clearpage

\begin{figure*}[p]\centering
\includegraphics[width=\textwidth,height=0.85\textheight,keepaspectratio]{figures/kudo_fig5_llama32-3b_L3_cot_v2.pdf}
\caption{Span-by-layer-window patching grids, three sources (the error-removal panel is empty: with CoT there are too few name errors to draw), after \citeauthor{kudo2026faithful}'s Figures~5--6.}\label{fig:dump-kudo-fig5-llama32-3b-L3-cot-v2}
\end{figure*}
\clearpage

\begin{figure*}[p]\centering
\includegraphics[width=\textwidth,height=0.85\textheight,keepaspectratio]{figures/kudo_fig5_llama32-3b_L3_direct_v2.pdf}
\caption{The same grids in the direct regime.}\label{fig:dump-kudo-fig5-llama32-3b-L3-direct-v2}
\end{figure*}
\clearpage

\begin{figure*}[p]\centering
\includegraphics[width=\textwidth,height=0.85\textheight,keepaspectratio]{figures/fig3_patching_llama32-3b_L3_v2.pdf}
\caption{Name errors removed by patching, by layer, with both controls.}\label{fig:dump-fig3-patching-llama32-3b-L3-v2}
\end{figure*}
\clearpage

\begin{figure*}[p]\centering
\includegraphics[width=\textwidth,height=0.85\textheight,keepaspectratio]{figures/kudo_fig2_llama32-3b_L3_cot_v1.pdf}
\caption{Accuracy heatmaps as above, chain-of-thought regime, queried variable.}\label{fig:dump-kudo-fig2-llama32-3b-L3-cot-v1}
\end{figure*}
\clearpage

\begin{figure*}[p]\centering
\includegraphics[width=\textwidth,height=0.85\textheight,keepaspectratio]{figures/kudo_fig2_llama32-3b_L3_direct_v1.pdf}
\caption{The same heatmaps in the direct regime, queried variable.}\label{fig:dump-kudo-fig2-llama32-3b-L3-direct-v1}
\end{figure*}
\clearpage

\begin{figure*}[p]\centering
\includegraphics[width=\textwidth,height=0.85\textheight,keepaspectratio]{figures/kudo_fig5_llama32-3b_L3_cot_v1.pdf}
\caption{Span-by-layer-window patching grids, chain-of-thought regime, queried variable; the error-removal panel is empty for the same reason.}\label{fig:dump-kudo-fig5-llama32-3b-L3-cot-v1}
\end{figure*}
\clearpage

\begin{figure*}[p]\centering
\includegraphics[width=\textwidth,height=0.85\textheight,keepaspectratio]{figures/kudo_fig5_llama32-3b_L3_direct_v1.pdf}
\caption{The same grids in the direct regime, queried variable.}\label{fig:dump-kudo-fig5-llama32-3b-L3-direct-v1}
\end{figure*}
\clearpage

\begin{figure*}[p]\centering
\includegraphics[width=\textwidth,height=0.85\textheight,keepaspectratio]{figures/fig3_patching_llama32-3b_L3_v1.pdf}
\caption{Name errors removed by patching, by layer, with both controls, queried variable.}\label{fig:dump-fig3-patching-llama32-3b-L3-v1}
\end{figure*}
\clearpage

\begin{figure*}[p]\centering
\includegraphics[width=\textwidth,height=0.85\textheight,keepaspectratio]{figures/fig_own_chain_llama32-3b_L3_cot_v1.pdf}
\caption{Llama-3.2-3B, queried variable: internal predictions during eight calculations the model generated before giving a wrong final answer.}\label{fig:dump-fig-own-chain-llama32-3b-L3-cot-v1}
\end{figure*}
\clearpage

\begin{figure*}[p]\centering
\includegraphics[width=\textwidth,height=0.85\textheight,keepaspectratio]{figures/fig3_patching_llama31-8b_L3_v2.pdf}
\caption{Llama-3.1-8B, name errors removed by patching, by layer, with both controls, direct regime, intermediate variable.}\label{fig:dump-fig3-patching-llama31-8b-L3-v2}
\end{figure*}
\clearpage

\begin{figure*}[p]\centering
\includegraphics[width=\textwidth,height=0.85\textheight,keepaspectratio]{figures/kudo_fig2_llama31-8b_L3_cot_v1.pdf}
\caption{Llama-3.1-8B, chain-of-thought regime, queried variable: the accuracy heatmaps of the Llama-3.2-3B figures. The value is decodable from the equation that restates the variable's definition, before the chain writes it.}\label{fig:dump-kudo-fig2-llama31-8b-L3-cot-v1}
\end{figure*}
\clearpage

\begin{figure*}[p]\centering
\includegraphics[width=\textwidth,height=0.85\textheight,keepaspectratio]{figures/kudo_fig2_llama31-8b_L3_cot_v2.pdf}
\caption{Llama-3.1-8B, chain-of-thought regime, intermediate variable.}\label{fig:dump-kudo-fig2-llama31-8b-L3-cot-v2}
\end{figure*}
\clearpage

\begin{figure*}[p]\centering
\includegraphics[width=\textwidth,height=0.85\textheight,keepaspectratio]{figures/kudo_fig2_llama31-8b_L3_direct_v2.pdf}
\caption{Llama-3.1-8B, direct regime, intermediate variable.}\label{fig:dump-kudo-fig2-llama31-8b-L3-direct-v2}
\end{figure*}
\clearpage

\begin{figure*}[p]\centering
\includegraphics[width=\textwidth,height=0.85\textheight,keepaspectratio]{figures/kudo_fig5_llama31-8b_L3_cot_v2.pdf}
\caption{Llama-3.1-8B, span-by-layer-window patching grids, chain-of-thought regime. The error-removal panel is empty: none of the 500 patched problems is answered with the name-suggested value.}\label{fig:dump-kudo-fig5-llama31-8b-L3-cot-v2}
\end{figure*}
\clearpage

\begin{figure*}[p]\centering
\includegraphics[width=\textwidth,height=0.85\textheight,keepaspectratio]{figures/kudo_fig5_llama31-8b_L3_direct_v2.pdf}
\caption{Llama-3.1-8B, the same grids in the direct regime.}\label{fig:dump-kudo-fig5-llama31-8b-L3-direct-v2}
\end{figure*}
\clearpage

\begin{figure*}[p]\centering
\includegraphics[width=\textwidth,height=0.85\textheight,keepaspectratio]{figures/kudo_fig2_llama31-8b_L3_direct_v1.pdf}
\caption{Llama-3.1-8B, direct regime, queried variable.}\label{fig:dump-kudo-fig2-llama31-8b-L3-direct-v1}
\end{figure*}
\clearpage

\begin{figure*}[p]\centering
\includegraphics[width=\textwidth,height=0.85\textheight,keepaspectratio]{figures/kudo_fig5_llama31-8b_L3_cot_v1.pdf}
\caption{Llama-3.1-8B, patching grids, chain-of-thought regime, queried variable; the error-removal panel is empty for the same reason.}\label{fig:dump-kudo-fig5-llama31-8b-L3-cot-v1}
\end{figure*}
\clearpage

\begin{figure*}[p]\centering
\includegraphics[width=\textwidth,height=0.85\textheight,keepaspectratio]{figures/kudo_fig5_llama31-8b_L3_direct_v1.pdf}
\caption{Llama-3.1-8B, patching grids, direct regime, queried variable.}\label{fig:dump-kudo-fig5-llama31-8b-L3-direct-v1}
\end{figure*}
\clearpage


% Generated from validated reviewer follow-up JSON.
\newsavebox{\roundfivebox}
\section{Task competence and value readouts}
\label{app:round5-competence}
Table~\ref{tab:round5-competence} compares both variables for 9 models. Layers are selected by neutral selectivity, averaging three probe seeds. Task accuracy and paired excess name errors pool demonstration sets 7, 11 and 13; probes use set 7. The within-model difference holds overall task accuracy fixed but does not isolate a causal effect of the readout.
We average the two variable readouts within each model before computing descriptive correlations. Table~\ref{tab:round5-correlation} repeats them after omitting each model. Roles are not independent observations. With this small, heterogeneous panel, correlations cannot distinguish overall competence from the role of decodable information.
\begin{table*}[t]
\centering\small
\sbox{\roundfivebox}{\begin{tabular}{@{}llrrrrr@{}}
\toprule
Model & Role & Layer & Task & Neutral probe & Misleading probe & Excess \\
\midrule
Gemma-3-12B-I & v1 & 40 & 100.0 & 100.0 & 100.0 & 0.0 \\
Gemma-3-12B-I & v2 & 42 & 100.0 & 100.0 & 100.0 & 0.0 \\
Gemma-3-4B & v1 & 25 & 100.0 & 100.0 & 100.0 & 0.0 \\
Gemma-3-4B & v2 & 28 & 100.0 & 99.8 & 100.0 & 0.0 \\
Gemma-3-4B-I & v1 & 29 & 100.0 & 100.0 & 99.9 & 0.0 \\
Gemma-3-4B-I & v2 & 28 & 100.0 & 99.8 & 99.9 & 0.0 \\
Llama-3.1-8B & v1 & 20 & 95.2 & 100.0 & 100.0 & -0.1 \\
Llama-3.1-8B & v2 & 22 & 95.2 & 99.4 & 99.2 & -0.0 \\
Llama-3.1-8B-I & v1 & 20 & 99.7 & 100.0 & 100.0 & 0.0 \\
Llama-3.1-8B-I & v2 & 20 & 99.7 & 100.0 & 100.0 & 0.0 \\
Llama-3.2-3B & v1 & 19 & 98.3 & 99.9 & 100.0 & 0.1 \\
Llama-3.2-3B & v2 & 21 & 98.3 & 97.6 & 99.6 & -0.2 \\
Llama-3.2-3B-I & v1 & 18 & 100.0 & 99.9 & 100.0 & 0.0 \\
Llama-3.2-3B-I & v2 & 19 & 100.0 & 99.9 & 99.9 & 0.0 \\
OLMo-2-1B & v1 & 14 & 47.2 & 99.7 & 100.0 & 0.4 \\
OLMo-2-1B & v2 & 0 & 47.2 & 21.8 & 21.8 & 3.2 \\
OLMo-2-7B & v1 & 22 & 95.3 & 100.0 & 100.0 & -0.1 \\
OLMo-2-7B & v2 & 26 & 95.3 & 99.4 & 97.6 & -0.0 \\
\bottomrule
\end{tabular}}
\ifdim\wd\roundfivebox>\textwidth\resizebox{\textwidth}{!}{\usebox{\roundfivebox}}\else\usebox{\roundfivebox}\fi
\caption{Level-3 arithmetic task and probe accuracies (\%) and paired excess name errors (points). Neutral task accuracy is shared by the two roles. All available complete model pairs are shown; below-floor models remain descriptive.}
\label{tab:round5-competence}
\end{table*}
\begin{table*}[t]
\centering\small
\sbox{\roundfivebox}{\begin{tabular}{@{}lrrr@{}}
\toprule
Omitted model & Task--excess & Probe--excess & Task--probe \\
\midrule
All models & -0.99 & -1.00 & +1.00 \\
Gemma-3-12B-I & -0.99 & -1.00 & +1.00 \\
Gemma-3-4B & -0.99 & -1.00 & +1.00 \\
Gemma-3-4B-I & -0.99 & -1.00 & +1.00 \\
Llama-3.1-8B & -0.99 & -1.00 & +1.00 \\
Llama-3.1-8B-I & -0.99 & -1.00 & +1.00 \\
Llama-3.2-3B & -0.99 & -1.00 & +1.00 \\
Llama-3.2-3B-I & -0.99 & -1.00 & +1.00 \\
OLMo-2-1B & +0.98 & +0.78 & +0.85 \\
OLMo-2-7B & -0.99 & -1.00 & +1.00 \\
\bottomrule
\end{tabular}}
\ifdim\wd\roundfivebox>\textwidth\resizebox{\textwidth}{!}{\usebox{\roundfivebox}}\else\usebox{\roundfivebox}\fi
\caption{Descriptive Pearson correlations of model means; Spearman values and within-model variable differences are in the saved summary. No significance or causal claim is based on these coefficients.}
\label{tab:round5-correlation}
\end{table*}
\begin{table*}[t]
\centering\small
\sbox{\roundfivebox}{\begin{tabular}{@{}lrrr@{}}
\toprule
Model & Neutral probe & Misleading probe & Excess \\
\midrule
Gemma-3-12B-I & +0.0 & +0.0 & +0.0 \\
Gemma-3-4B & -0.2 & +0.0 & +0.0 \\
Gemma-3-4B-I & -0.2 & +0.0 & +0.0 \\
Llama-3.1-8B & -0.6 & -0.8 & +0.1 \\
Llama-3.1-8B-I & +0.0 & +0.0 & +0.0 \\
Llama-3.2-3B & -2.2 & -0.4 & -0.3 \\
Llama-3.2-3B-I & +0.1 & -0.1 & +0.0 \\
OLMo-2-1B & -77.9 & -78.2 & +2.7 \\
OLMo-2-7B & -0.5 & -2.3 & +0.1 \\
\bottomrule
\end{tabular}}
\ifdim\wd\roundfivebox>\textwidth\resizebox{\textwidth}{!}{\usebox{\roundfivebox}}\else\usebox{\roundfivebox}\fi
\caption{Intermediate minus queried variable, in percentage points. Overall neutral task accuracy is identical within each model. These differences are descriptive comparisons between variable positions.}
\label{tab:round5-within}
\end{table*}
\section{Gold-trained transfer to generated chains}
\label{app:round5-generated}
We read the state before the first numeric assignment to the target in each saved generated chain. Expression operands are not counted as completed value writes; tokens that include any part of the supplied value are excluded. A fixed sample of 200 matched sets is selected by a hash independent of outcomes, shared across models, variables and demonstration sets. All additional chains with name errors and their neutral twins are evaluated separately.
Neutral gold-chain probes use 10,000 separate training problems, three probe seeds and the original neutral-selectivity layer. Table~\ref{tab:round5-generated} tests transfer to held-out neutral generated chains. A high-accuracy interpretation requires at least 90\% neutral calibration in that demonstration set. Cases below that threshold yield descriptive readouts only. The augmentation is selected by observed errors and cannot estimate their population rate. Appendix~\ref{app:round6-generated} separately trains probes on generated neutral chains and tests the same held-out prefixes.
Llama-3.2-3B meets the neutral transfer threshold in 6 of 6 cells, with calibration accuracy from 97.2 to 100.0\%.
OLMo-2-1B meets the neutral transfer threshold in 0 of 6 cells, with calibration accuracy from 19.4 to 62.2\%.
Name-error cases in the competent model are few; error-conditioned readouts cannot establish a general mechanism from that subset. Failed neutral calibration prevents a strong generated-chain interpretation for the exception model.
\begin{table*}[t]
\centering\small
\sbox{\roundfivebox}{\begin{tabular}{@{}lrrlllll@{}}
\toprule
Model & Set & Role & Neutral coverage & Calibration & Sample error probe & Extra errors $n$ & Extra error probe \\
\midrule
OLMo-2-1B & 7 & v1 & 194/200 & 62.2 [55.2, 68.7] & 0.0 [0.0, 0.0] & 69 & 1.4 [0.0, 3.9] \\
OLMo-2-1B & 7 & v2 & 186/200 & 19.4 [14.0, 25.3] & 54.5 [27.3, 81.8] & 82 & 65.9 [56.1, 75.6] \\
OLMo-2-1B & 11 & v1 & 187/200 & 55.3 [48.3, 62.6] & 0.0 [0.0, 0.0] & 59 & 5.6 [0.6, 12.4] \\
OLMo-2-1B & 11 & v2 & 129/200 & 21.7 [14.7, 28.7] & 100.0 [100.0, 100.0] & 21 & 66.7 [47.6, 85.7] \\
OLMo-2-1B & 13 & v1 & 193/200 & 52.7 [45.6, 59.8] & 27.5 [7.8, 49.0] & 159 & 10.7 [6.5, 15.3] \\
OLMo-2-1B & 13 & v2 & 124/200 & 24.2 [16.9, 32.3] & 100.0 [100.0, 100.0] & 30 & 60.0 [43.3, 76.7] \\
Llama-3.2-3B & 7 & v1 & 200/200 & 98.0 [96.0, 99.5] & 0.0 [0.0, 0.0] & 8 & 12.5 [0.0, 37.5] \\
Llama-3.2-3B & 7 & v2 & 200/200 & 97.2 [94.8, 99.0] & -- & 0 & -- \\
Llama-3.2-3B & 11 & v1 & 199/200 & 100.0 [100.0, 100.0] & -- & 0 & -- \\
Llama-3.2-3B & 11 & v2 & 200/200 & 100.0 [100.0, 100.0] & -- & 0 & -- \\
Llama-3.2-3B & 13 & v1 & 199/200 & 99.5 [98.5, 100.0] & -- & 5 & 0.0 [0.0, 0.0] \\
Llama-3.2-3B & 13 & v2 & 200/200 & 99.5 [98.5, 100.0] & -- & 0 & -- \\
\bottomrule
\end{tabular}}
\ifdim\wd\roundfivebox>\textwidth\resizebox{\textwidth}{!}{\usebox{\roundfivebox}}\else\usebox{\roundfivebox}\fi
\caption{Generated-chain true-digit accuracy (\%) [95\% interval]. Sample error probe conditions on name errors in the fixed sample; extra name-error cases exclude that sample. Dashes indicate no eligible cases. Saved summaries record all missing or excluded boundaries and calibration decisions.}
\label{tab:round5-generated}
\end{table*}

% Generated from validated round-six summary.
\newsavebox{\roundsixbox}
\section{Training probes on generated neutral chains}
\label{app:round6-generated}
\input{figures/readout_check}
Gold-trained probe transfer (Appendix~\ref{app:round5-generated}) may fail because the states along generated chains differ from those along the supplied correct chain. We therefore train on freely generated neutral chains from \NUM{r6-source-train_n} source problems and select layers on \NUM{r6-source-validation_n} separate validation problems. Canonical equations partition the source pool, keeping renamed versions of a computation together. The training and validation pools contain \NUM{r6-source-train_keys} and \NUM{r6-source-validation_keys} computations; both exclude all test and demonstration computations. These counts distinguish independent computations from repeated name renderings.
Generation is greedy with demonstration set 7 and a 128-token limit. The probe sees the hidden state before the target variable's first numeric assignment, with every token boundary checked to exclude the value digit. Missing or unsafe boundaries are excluded and counted. We retain all usable neutral chains regardless of whether the model writes the correct value, using the executed true value as the training label. Three probe seeds use the original full-batch SGD recipe (learning rate 0.001, 10,000 epochs). States are stored in float32 and nonfinite states or fitted weights are rejected. Layers are searched at stride two, including the final layer, and chosen only by mean validation accuracy with the lowest-layer tie break. The name-identity control is trained separately and does not select layers.
Table~\ref{tab:round6-selection} reports the usable source counts and selected layers. Evaluation uses exactly the earlier fixed 200-set sample and separately reported additional name-error cases under demonstration sets 7, 11 and 13. Table~\ref{tab:round6-calibration} gives held-out neutral calibration and boundary coverage. Intervals use 4,000 canonical-computation bootstrap draws, keeping repeated names together. The 90\% calibration threshold applies to the point estimate, not the lower interval endpoint.
Llama passes all six calibration cells; OLMo passes none. Retraining on generated chains therefore does not validate a strong readout account of OLMo's own errors. Table~\ref{tab:round6-errors} retains its error readouts descriptively. Llama has only one name-error case in the primary sample and thirteen additional observed writes; these rare, selected cases cannot establish a general mechanism. Paired changes from gold-trained transfer are shown in Table~\ref{tab:round6-changes}; they compare the same eligible prefixes, not independent model samples.
\begin{table*}[t]
\centering\small
\sbox{\roundsixbox}{\begin{tabular}{@{}llrrrrr@{}}
\toprule
Model & Role & Train $n$ & Validation $n$ & Layer & Validation (\%) & Name ID (\%) \\
\midrule
OLMo-2-1B-I & v1 & 1924 & 390 & 14 & 47.9 & 41.6 \\
OLMo-2-1B-I & v2 & 1877 & 376 & 16 & 87.8 & 86.6 \\
Llama-3.2-3B & v1 & 1994 & 398 & 22 & 97.5 & 26.8 \\
Llama-3.2-3B & v2 & 2000 & 400 & 20 & 99.2 & 28.6 \\
\bottomrule
\end{tabular}}
\ifdim\wd\roundsixbox>\textwidth\resizebox{\textwidth}{!}{\usebox{\roundsixbox}}\else\usebox{\roundsixbox}\fi
\caption{Probe source rows with usable generated boundaries, selected layers and validation scores. Source cohorts are computation-disjoint; usable counts vary by model and variable. Validation scores select layers and are not test estimates.}
\label{tab:round6-selection}
\end{table*}
\begin{table*}[t]
\centering\small
\sbox{\roundsixbox}{\begin{tabular}{@{}llrlllc@{}}
\toprule
Model & Role & Set & Neutral coverage & Calibration (\%) & Name ID (\%) & Pass \\
\midrule
OLMo-2-1B-I & v1 & 7 & 194/200 & 61.3 [54.1, 68.6] & 37.3 [31.2, 43.4] & no \\
OLMo-2-1B-I & v1 & 11 & 187/200 & 52.8 [45.1, 60.3] & 32.4 [26.6, 38.5] & no \\
OLMo-2-1B-I & v1 & 13 & 193/200 & 50.3 [43.0, 57.4] & 35.2 [29.1, 41.6] & no \\
OLMo-2-1B-I & v2 & 7 & 186/200 & 78.0 [71.5, 83.9] & 81.9 [76.9, 86.7] & no \\
OLMo-2-1B-I & v2 & 11 & 129/200 & 80.1 [72.8, 86.8] & 56.6 [49.6, 63.3] & no \\
OLMo-2-1B-I & v2 & 13 & 124/200 & 71.8 [64.2, 78.4] & 46.8 [37.8, 55.6] & no \\
Llama-3.2-3B & v1 & 7 & 200/200 & 98.5 [96.7, 100.0] & 18.5 [13.8, 23.7] & yes \\
Llama-3.2-3B & v1 & 11 & 199/200 & 100.0 [100.0, 100.0] & 13.9 [9.5, 18.5] & yes \\
Llama-3.2-3B & v1 & 13 & 199/200 & 97.2 [94.2, 99.5] & 16.9 [12.3, 21.8] & yes \\
Llama-3.2-3B & v2 & 7 & 200/200 & 96.3 [93.6, 98.6] & 28.8 [22.7, 35.0] & yes \\
Llama-3.2-3B & v2 & 11 & 200/200 & 100.0 [100.0, 100.0] & 17.2 [12.4, 22.1] & yes \\
Llama-3.2-3B & v2 & 13 & 200/200 & 97.7 [95.5, 99.3] & 17.8 [13.3, 22.8] & yes \\
\bottomrule
\end{tabular}}
\ifdim\wd\roundsixbox>\textwidth\resizebox{\textwidth}{!}{\usebox{\roundsixbox}}\else\usebox{\roundsixbox}\fi
\caption{Held-out primary neutral generated chains. Coverage counts valid pre-value boundaries over all primary neutral rows. Accuracy intervals are 95\% canonical-computation bootstrap intervals; pass means the calibration point estimate meets 90\%.}
\label{tab:round6-calibration}
\end{table*}
\begin{table*}[t]
\centering\small
\sbox{\roundsixbox}{\begin{tabular}{@{}llrrlrl@{}}
\toprule
Model & Role & Set & Sample $n$ & Sample accuracy & Extra $n$ & Extra accuracy \\
\midrule
OLMo-2-1B-I & v1 & 7 & 8 & 4.2 [0.0, 12.5] & 69 & 2.9 [0.0, 7.0] \\
OLMo-2-1B-I & v1 & 11 & 7 & 19.0 [0.0, 38.1] & 59 & 10.7 [4.2, 18.8] \\
OLMo-2-1B-I & v1 & 13 & 17 & 21.6 [7.8, 39.2] & 159 & 13.6 [8.6, 19.4] \\
OLMo-2-1B-I & v2 & 7 & 11 & 33.3 [15.2, 51.5] & 82 & 18.7 [11.4, 27.2] \\
OLMo-2-1B-I & v2 & 11 & 2 & 0.0 [0.0, 0.0] & 21 & 17.5 [3.3, 31.8] \\
OLMo-2-1B-I & v2 & 13 & 1 & 0.0 [0.0, 0.0] & 30 & 25.6 [13.7, 39.1] \\
Llama-3.2-3B & v1 & 7 & 1 & 0.0 [0.0, 0.0] & 8 & 12.5 [0.0, 42.9] \\
Llama-3.2-3B & v1 & 11 & 0 & -- & 0 & -- \\
Llama-3.2-3B & v1 & 13 & 0 & -- & 5 & 0.0 [0.0, 0.0] \\
Llama-3.2-3B & v2 & 7 & 0 & -- & 0 & -- \\
Llama-3.2-3B & v2 & 11 & 0 & -- & 0 & -- \\
Llama-3.2-3B & v2 & 13 & 0 & -- & 0 & -- \\
\bottomrule
\end{tabular}}
\ifdim\wd\roundsixbox>\textwidth\resizebox{\textwidth}{!}{\usebox{\roundsixbox}}\else\usebox{\roundsixbox}\fi
\caption{True-digit accuracy (\%) [95\% interval] on actual name errors, with the fixed sample and additional error-selected cases separated. Extra cases exclude the primary sample. Dashes denote empty strata. OLMo fails neutral calibration, so its readouts are descriptive.}
\label{tab:round6-errors}
\end{table*}
\begin{table*}[t]
\centering\small
\sbox{\roundsixbox}{\begin{tabular}{@{}llrlll@{}}
\toprule
Model & Role & Set & Neutral change & Sample error change & Extra error change \\
\midrule
OLMo-2-1B-I & v1 & 7 & -0.9 [-4.9, 3.4] & 4.2 [0.0, 12.5] & 1.4 [-1.0, 5.3] \\
OLMo-2-1B-I & v1 & 11 & -2.5 [-7.7, 2.7] & 19.0 [0.0, 38.1] & 5.1 [-0.6, 12.1] \\
OLMo-2-1B-I & v1 & 13 & -2.4 [-7.2, 2.2] & -5.9 [-19.6, 5.9] & 2.9 [-0.8, 6.7] \\
OLMo-2-1B-I & v2 & 7 & 58.6 [49.5, 67.5] & -21.2 [-57.6, 12.2] & -47.2 [-63.3, -30.1] \\
OLMo-2-1B-I & v2 & 11 & 58.4 [47.9, 67.9] & -100.0 [-100.0, -100.0] & -49.2 [-71.4, -20.8] \\
OLMo-2-1B-I & v2 & 13 & 47.6 [37.7, 57.1] & -100.0 [-100.0, -100.0] & -34.4 [-55.6, -13.3] \\
Llama-3.2-3B & v1 & 7 & 0.5 [0.0, 1.6] & 0.0 [0.0, 0.0] & 0.0 [0.0, 0.0] \\
Llama-3.2-3B & v1 & 11 & 0.0 [0.0, 0.0] & -- & -- \\
Llama-3.2-3B & v1 & 13 & -2.3 [-5.1, -0.3] & -- & 0.0 [0.0, 0.0] \\
Llama-3.2-3B & v2 & 7 & -0.8 [-2.9, 0.9] & -- & -- \\
Llama-3.2-3B & v2 & 11 & 0.0 [0.0, 0.0] & -- & -- \\
Llama-3.2-3B & v2 & 13 & -1.8 [-3.7, -0.5] & -- & -- \\
\bottomrule
\end{tabular}}
\ifdim\wd\roundsixbox>\textwidth\resizebox{\textwidth}{!}{\usebox{\roundsixbox}}\else\usebox{\roundsixbox}\fi
\caption{Paired accuracy changes in percentage points [95\% interval] from gold-trained transfer to generated-neutral training on the same valid prefixes. Error-selection and failed neutral calibration limit the interpretation.}
\label{tab:round6-changes}
\end{table*}


\end{document}
